\section*{Capítulo \ref{ch:b2:magmatic-differentiation} --- A diferenciação magmática}

\begin{solution}{exo:b2:magmatic-differentiation:1}
(a) O liquidus de An$_{50}$ é \qty{1450}{\celsius}. O primeiro cristal fica sobre o
solidus a essa temperatura, entre An$_{66.7}$ (1372) e An$_{83.3}$ (1465):
$66.7 + 16.6\times(1450 - 1372)/(1465 - 1372) = 66.7 + 13.9 = \qty{80.6}{\percent}$
de anortita. (b) A \qty{1394}{\celsius}, o líquido está sobre o liquidus,
An$_{33.3}$; os cristais, sobre o solidus, $66.7 + 16.6\times22/93 =
\qty{70.6}{\percent}$. Pela regra da alavanca, a fração cristalizada é $(50 -
33.3)/(70.6 - 33.3) = 16.7/37.3 = 0.448$. (c) O último líquido desaparece quando os
cristais chegam a An$_{50}$, à temperatura do solidus \qty{1287}{\celsius}; o
liquidus ali fica entre An$_{11.1}$ (1265) e An$_{20}$ (1334): $11.1 +
8.9\times22/69 = \qty{13.9}{\percent}$ de anortita.
\end{solution}

\begin{solution}{exo:b2:magmatic-differentiation:2}
$\bar D_{\ce{Sr}} = 0.15\times0.000063 + 0.35\times0.157 + 0.50\times2.7 = 0.00001 +
0.0550 + 1.35 = 1.40$; $\bar D_{\ce{Zr}} = 0.0001 + 0.0683 + 0.0005 = 0.069$;
$\bar D_{\ce{Rb}} = 0.00001 + 0.0012 + 0.0125 = 0.014$. Só o estrôncio, retido pelo
plagioclásio, é compatível ($\bar D > 1$); o zircônio e o rubídio são incompatíveis,
o rubídio ainda mais.
\end{solution}

\begin{solution}{exo:b2:magmatic-differentiation:3}
(a) $C_L/C_0 = 0.5^{D-1}$: Sr $0.5^{0.40} = 0.758$ (ele cai); Zr $0.5^{-0.931} =
1.907$; Rb $0.5^{-0.986} = 1.981$, perto do limite 2. (b) Zr/Rb varia por um fator de
$0.5^{0.069 - 0.014} = 0.5^{0.055} = 0.963$, menos de \qty{4}{\percent}, enquanto cada
elemento quase dobrou: como diz o corolário, a razão de dois \omterm{def:b2:elements-in-the-earth:partition}{elementos
incompatíveis} conserva o valor do \omterm{def:b2:magma-emplacement:magma}{magma} parental, e assim informa sobre a fonte, não
sobre a cristalização.
\end{solution}

\begin{solution}{exo:b2:magmatic-differentiation:4}
(a) \ce{SiO2} $+5.03$, \ce{MgO} $-1.95$, \ce{K2O} $+0.82$ (por cento), a potassa
multiplicada por $1.67/0.85 = 1.96$. (b) Para um elemento perfeitamente incompatível,
$C_L/C_0 = 1/F$, logo $F = 0.85/1.67 = 0.509$: metade do \omterm{def:b2:magma-emplacement:magma}{magma} teria de
cristalizar. (c) Remover \qty{12.5}{g} por \qty{100}{g} deixa $F \approx 0.88$ e
elevaria a potassa apenas por um fator de cerca de $1/0.88 = 1.14$: a cristalização sozinha não pode
explicar o aumento. A potassa veio de fora, dos xenólitos graníticos, ricos em
potassa (\qty{3.6}{\percent}), que o \omterm{def:b2:magma-emplacement:magma}{magma} assimilou.
\end{solution}

\begin{solution}{exo:b2:magmatic-differentiation:5}
(a) $F = 0.8$, $D = 5$. Fracionada: $0.8^{4} = 0.410$; equilíbrio:
$1/(5 + 0.8\times(1 - 5)) = 1/1.8 = 0.556$. (b) $F = 0.1$, $D = 0.01$. Fracionada:
$0.1^{-0.99} = 9.77$; equilíbrio: $1/(0.01 + 0.1\times0.99) = 1/0.109 = 9.17$.
(c) Quando $F \to 0$, a lei de Rayleigh tende a $0$ para $D > 1$ e ao infinito para
$D < 1$, ao passo que a lei de equilíbrio tende a $1/D$ (o último líquido em equilíbrio
com cristais da composição global). Só a \omterm{def:b2:magmatic-differentiation:crystallisation}{cristalização fracionada} pode esgotar completamente um
líquido de um \omterm{def:b2:elements-in-the-earth:partition}{elemento compatível}.
\end{solution}

\begin{solution}{exo:b2:magmatic-differentiation:6}
(a) Uma unidade de fórmula contém 1.6 \ce{MgO}, 0.4 \ce{FeO} e um \ce{SiO2}: $1.6\times
40.30 = 64.48$, $0.4\times71.84 = 28.74$, 60.08, total \qty{153.30}{g}. Porcentagens
em massa: \ce{MgO} \qty{42.06}{\percent}, \ce{FeO} \qty{18.75}{\percent},
\ce{SiO2} \qty{39.19}{\percent}. (b) Regra da alavanca sobre \ce{MgO}: $x = (10.0 - 8.0)/(42.06
- 8.0) = 2.0/34.06 = 0.0587$, cerca de \qty{6}{\percent} de olivina. (c) Sílica do
líquido: $(49.0 - 0.0587\times39.19)/(1 - 0.0587) = 46.70/0.9413 =
\qty{49.61}{\percent}$. Remover olivina, mais pobre em sílica que o basalto, eleva a
sílica do líquido, mas só lentamente.
\end{solution}

\begin{solution}{exo:b2:magmatic-differentiation:7}
(a) $F(0) = 0.768$; $F(1000) = 0.1091 - 0.5906 + 0.7678 = 0.286$; $F(2000) = 0.4364
- 1.1813 + 0.7678 = 0.023$. (b) O discriminante é $(5.9064\times10^{-4})^2 -
4\times1.091\times10^{-7}\times0.7678 = 3.4885\times10^{-7} - 3.3507\times10^{-7} =
1.379\times10^{-8}$, de raiz quadrada $1.174\times10^{-4}$; a raiz menor é
$H = (5.9064 - 1.174)\times10^{-4}/(2.182\times10^{-7}) = \qty{2169}{m}$, o
Horizonte Sanduíche. (c) $1 - 0.768 = 0.232$: quase um quarto do \omterm{def:b2:magma-emplacement:magma}{magma} tinha
cristalizado antes de se formar a base da Zona Inferior exposta, na Zona Oculta
abaixo dela e na Série Marginal ao longo das paredes e do teto.
\end{solution}

\begin{solution}{exo:b2:magmatic-differentiation:8}
O \omterm{def:b2:magmatic-differentiation:cumulate}{cumulato} contém $D\,C_L = D\,C_0F^{D-1}$, de modo que, entre dois níveis,
$\ln(\mathrm{Ni}_2/\mathrm{Ni}_1) = (D - 1)\ln(F_2/F_1)$. (a) $D - 1 = \ln(54/177)/
\ln(0.29/0.70) = (-1.187)/(-0.881) = 1.35$, logo $D \approx 2.35$ (a regressão sobre
os 66 \omterm{def:b2:magmatic-differentiation:cumulate}{cumulatos} com $F \ge 0.23$ dá 2.38). (b) $D - 1 = \ln(6.2/54)/\ln(0.17/0.29) = (-2.164)/(-0.534)
= 4.05$, logo $D \approx 5.05$. (c) O $D$ global muda quando a assembleia que
cristaliza muda (um novo mineral que capta níquel, ou um que sai), e os
próprios coeficientes mineral--fundido aumentam à medida que o líquido se resfria e muda de
composição, como foi medido para a olivina; o líquido aprisionado nos \omterm{def:b2:magmatic-differentiation:cumulate}{cumulatos} borra a
estimativa.
\end{solution}

\begin{solution}{exo:b2:magmatic-differentiation:9}
(a) $100 + 25.4 - 2.9 - 9.6 = \qty{112.9}{g}$. (b) Gramas de \ce{SiO2}: $55.53 +
25.4\times0.7248 - 2.9\times0.3919 - 9.6\times0.5054 = 55.53 + 18.41 - 1.14 - 4.85 =
\qty{67.95}{g}$, logo $67.95/112.9 = \qty{60.19}{\percent}$, contra 60.16. \ce{MgO}:
$4.66 + 0.20 - 1.22 = \qty{3.64}{g}$, \qty{3.22}{\percent} contra 3.19. \ce{K2O}:
$0.88 + 0.92 = \qty{1.80}{g}$, \qty{1.60}{\percent} contra 1.75. (c) A sílica e a magnésia, dois dos três
controles, ajustam-se a menos de 0.03; a potassa, não usada como controle, fica baixa em 0.15, quase um
décimo de seu valor: o material assimilado tinha mais potassa que o xenólito
médio, ou um fundido rico em potassa foi tomado preferencialmente.
\end{solution}

\begin{solution}{exo:b2:magmatic-differentiation:10}
(a) $z = (0.2 + 1.5 - 1)/(0.2 - 1) = 0.7/(-0.8) = -0.875$; $F^{-z} = 0.7^{0.875} =
0.7319$; $C^*/C_0 = 0.2\times0.5/(0.2 - 1 + 1.5) = 0.1/0.7 = 0.1429$, logo $C/C_0 =
0.7319 + 0.1429\times(1 - 0.7319) = 0.7319 + 0.0383 = 0.770$. (b) Rayleigh:
$0.7^{0.5} = 0.837$. O \omterm{def:b2:magma-emplacement:magma}{magma} AFC tem menos estrôncio por duas razões: para perder
\qty{30}{\percent} de sua massa enquanto ganha $r$ vezes o que cristaliza, ele precisa
cristalizar $(1 - 0.7)/(1 - 0.2) = 0.375$ em vez de 0.30, e o estrôncio que ganha
vem de uma rocha duas vezes mais pobre. (c) $(0.7319\times0.7035 + 0.0383\times0.7200)/0.7702
= (0.5149 + 0.0276)/0.7702 = 0.70432$: a razão subiu 0.0008, ao passo que a
cristalização sozinha a teria deixado em 0.7035.
\end{solution}

\begin{solution}{exo:b2:magmatic-differentiation:11}
(a) $1.1\times(1100 - 200) + 250 = 990 + 250 = \qty{1240}{J}$ por grama. (b)
$25.4\times1240 = \qty{31496}{J}$, perto dos \qty{31350}{J} de Wilcox. (c)
$31496/370 = \qty{85.1}{g}$ de plagioclásio. (d) O balanço de massa removia apenas $2.9 +
9.6 = \qty{12.5}{g}$ de cristais, cujo calor latente (cerca de \qty{4600}{J}) cobre menos
de um sexto da necessidade; $r = 25.4/12.5 = 2.03$, duas vezes mais rocha assimilada que
cristais formados, ao passo que o balanço térmico admite $r$ bem abaixo de 1. O \omterm{def:b2:magma-emplacement:magma}{magma} emitido
não pode ter fornecido o calor por si mesmo: uma massa maior de \omterm{def:b2:magma-emplacement:magma}{magma} deve ter cristalizado
e se resfriado em profundidade, longe da vista.
\end{solution}

\begin{solution}{exo:b2:magmatic-differentiation:12}
(a) Etapa 1, $F$ de 1 a 0.8 com $D = 5$: fator $0.8^{4} = 0.4096$. Etapa 2,
de 0.8 a 0.5 com $D = 1.5$: a lei se aplica ao líquido restante, de modo que o
fator é $(0.5/0.8)^{0.5} = 0.7906$. Níquel: $250\times0.4096\times0.7906 =
\qty{81.0}{ppm}$. (b) $250\times0.5^{4} = \qty{15.6}{ppm}$, cinco vezes baixo demais. (c)
A lei de Rayleigh supõe um $D$ constante; quando a assembleia que cristaliza muda,
$D$ salta, e a lei deve ser reiniciada a partir da composição do líquido atingida na
mudança, com $F$ medido em relação a esse líquido.
\end{solution}

\begin{solution}{pb:b2:magmatic-differentiation:1}
\noindent\textbf{Parte I.}

\textbf{1.} $\bar D_{\ce{Ni}} = 0.20\times5.0 + 0.50\times0 + 0.30\times2.0 = 1.60$.

\textbf{2.} $\bar D_{\ce{Zr}} = 0.20\times0.0007 + 0.50\times0.001 + 0.30\times0.195 =
0.00014 + 0.0005 + 0.0585 = 0.059$.

\textbf{3.} Zona superior: $\bar D_{\ce{Ni}} = 0.45\times2.0 = 0.90$; $\bar D_{\ce{Zr}} =
0.55\times0.001 + 0.45\times0.195 = 0.088$.

\textbf{4.} O níquel é compatível na zona inferior (1.60) e se torna ligeiramente
incompatível (0.90) quando a olivina desaparece; o zircônio é incompatível nas duas.

\textbf{5.} O clinopiroxênio: ele fornece $0.0585$ dos $0.059$ da zona inferior, quase
tudo, porque é o único mineral com um coeficiente longe de zero.

\textbf{6.} $(C_{\ce{Ni}}/C_{\ce{Zr}})$ varia por um fator de $F^{D_{\ce{Ni}} - D_{\ce{Zr}}}$;
a cada redução à metade, $0.5^{1.60 - 0.059} = 0.5^{1.541} = 0.344$.

\noindent\textbf{Parte II.}

\textbf{7.} $1.60\times150 = \qty{240}{ppm}$.

\textbf{8.} O líquido tinha $190/1.60 = \qty{118.75}{ppm}$; $F = (118.75/150)^{1/0.6} =
0.7917^{1.667} = 0.68$.

\textbf{9.} $80\times0.68^{-0.941} = 80\times1.438 = \qty{115}{ppm}$.

\textbf{10.} Fracionada: $150\times0.68^{0.6} = 150\times0.793 = \qty{119.0}{ppm}$;
equilíbrio: $150/(1.6 + 0.68\times(1 - 1.6)) = 150\times0.839 = \qty{125.8}{ppm}$.
A diferença é de apenas \qty{6}{\percent} depois que um terço do \omterm{def:b2:magma-emplacement:magma}{magma}
cristalizou: para um elemento moderadamente compatível, as duas leis permanecem próximas no início.

\textbf{11.} $150\times0.55^{0.6} = \qty{104.8}{ppm}$.

\textbf{12.} Logo abaixo: $1.60\times104.8 = \qty{167.7}{ppm}$; logo acima:
$0.90\times104.8 = \qty{94.3}{ppm}$. O perfil de níquel dos \omterm{def:b2:magmatic-differentiation:cumulate}{cumulatos} cai
abruptamente \qty{44}{\percent} no nível em que a olivina desaparece das
camadas: um degrau, não uma curva suave.

\textbf{13.} Líquido: $104.8\times(0.40/0.55)^{0.90 - 1} = 104.8\times1.032 =
\qty{108.2}{ppm}$; \omterm{def:b2:magmatic-differentiation:cumulate}{cumulato} $0.90\times108.2 = \qty{97.4}{ppm}$. Com $\bar D < 1$,
o níquel agora sobe lentamente no líquido e nos \omterm{def:b2:magmatic-differentiation:cumulate}{cumulatos}: o perfil, depois de seu
degrau, sobe suavemente.

\noindent\textbf{Parte III.}

\textbf{14.} $x_{\mathrm{ol}}(m_{\mathrm{ol},j} - d_j) + x_{\mathrm{pl}}
(m_{\mathrm{pl},j} - d_j) + x_{\mathrm{cpx}}(m_{\mathrm{cpx},j} - d_j) = 100\,(p_j -
d_j)$.

\textbf{15.} \ce{MgO}: $37.16\,x_{\mathrm{ol}} - 4.9\,x_{\mathrm{pl}} +
11.1\,x_{\mathrm{cpx}} = 410$; \ce{CaO}: $-7.3\,x_{\mathrm{ol}} +
7.06\,x_{\mathrm{pl}} + 13.2\,x_{\mathrm{cpx}} = 420$; \ce{Al2O3}:
$-14.1\,x_{\mathrm{ol}} + 17.6\,x_{\mathrm{pl}} - 11.1\,x_{\mathrm{cpx}} = 190$.

\textbf{16.} Resolvendo (por eliminação ou pela regra de Cramer): $x_{\mathrm{ol}} = 9.03$,
$x_{\mathrm{pl}} = 30.84$, $x_{\mathrm{cpx}} = 20.32$ gramas por \qty{100}{g}.

\textbf{17.} Cristais $9.03 + 30.84 + 20.32 = \qty{60.19}{g}$; líquido
\qty{39.81}{g}: $F = 0.398$.

\textbf{18.} $(4950 - 9.03\times39.19 - 30.84\times50.54 - 20.32\times51.5)/39.81 =
(4950 - 353.9 - 1558.7 - 1046.5)/39.81 = 1990.9/39.81 = \qty{50.0}{\percent}$, como no
\omterm{def:b2:magma-emplacement:bodies}{dique}: o óxido não usado como controle é reproduzido.

\textbf{19.} Os cristais removidos são $9.03/60.19 = \qty{15}{\percent}$ de olivina,
\qty{51}{\percent} de plagioclásio e \qty{34}{\percent} de clinopiroxênio, entre as
modas das duas zonas. Eles correspondem à assembleia inferior cristalizando de $F = 1$ a
0.55 ($0.45\times20 = 9.0$, $0.45\times50 = 22.5$, $0.45\times30 = 13.5$) seguida da
superior de 0.55 a 0.40 ($0.15\times55 = 8.25$, $0.15\times45 = 6.75$): 9.0,
30.75 e 20.25 gramas.

\textbf{20.} $0.30/0.75 = 0.40$.

\noindent\textbf{Parte IV.}

\textbf{21.} $0.55^{0.059 - 1} = 0.55^{-0.941} = 1.755$; $80\times1.755 =
\qty{140.4}{ppm}$.

\textbf{22.} $188/140.4 = 1.339 = (F/0.55)^{0.088 - 1}$, logo $F = 0.55\times1.339^{-1/
0.912} = 0.55\times0.726 = 0.399$.

\textbf{23.} $(188/80)^{1/(0.07 - 1)} = 2.35^{-1.075} = 0.399$. Para \omterm{def:b2:elements-in-the-earth:partition}{elementos
incompatíveis}, $C_L/C_0 \approx 1/F$ qualquer que seja o $D$ pequeno: a resposta quase não depende
dos detalhes de $D$, o que faz dos \omterm{def:b2:elements-in-the-earth:partition}{elementos incompatíveis} os melhores medidores de $F$.

\textbf{24.} O plagioclásio, que não contém ferro, forma metade dos cristais, e o
piroxênio e a olivina juntos contêm menos ferro do que magnésia: o líquido
se enriquece em ferro e se empobrece em magnésia sem ganhar sílica. Esse enriquecimento
em ferro com sílica quase constante é a marca da série toleítica do
capítulo seguinte.

\textbf{25.} Cristais cúmulus puros conteriam $0.088\times188 = \qty{16.5}{ppm}$. Com uma
fração em massa $t$ de líquido aprisionado, $16.5(1 - t) + 188\,t = 47$, logo $t = (47 -
16.5)/(188 - 16.5) = 30.5/171.5 = 0.18$: cerca de \qty{18}{\percent} de líquido aprisionado.

\textbf{26.} As três estimativas, elementos maiores (0.398), potassa (0.40) e
zircônio (0.399), concordam: quando a camada mais alta cristalizou,
\textbf{\qty[text-series-to-math=true]{40}{\percent}} do \omterm{def:b2:magma-emplacement:magma}{magma} original ainda estava
líquido.
\end{solution}
